\chapter{证券担保品估值及价差风险计量技术研究进展综述}

证券担保品估值及价差风险计量涉及资产定价、金融计量、概率预测、自然语言处理和模型风险管理等多个研究领域。现有研究已经分别形成了股票波动率建模、债券现金流折现、基金净值跟踪、期货期限结构、期权无套利定价、信用风险强度模型以及 VaR/ES 返回检验等较为成熟的技术路线，但这些路线通常围绕单一产品、单一目标或单一评价口径展开。对于同时覆盖六类证券担保品的科技报告，研究重点不只是汇集若干算法，而是回答三个相互衔接的问题：如何在保留产品定价约束的前提下形成统一的持有期价值与损失输出，如何检验文本和大模型是否提供了可识别的增量信息，以及如何用真实市场损益、严格时间切分和持续监控验证上述输出。本章按照这三条主线综述代表性研究，并据此说明本课题的技术切入点。

\section{多品种证券担保品估值与风险计量技术的研究进展}

多品种证券担保品估值技术的发展总体经历了“确定性现金流或无套利定价—随机波动与条件分布—跨序列概率预测—风险与流动性联合刻画”的演进过程。经典方法的优势在于经济含义清晰、参数与产品条款可以对应；现代统计学习方法的优势在于能够共享大量资产之间的信息并刻画非线性关系。两类方法并非简单替代关系：如果模型只追求样本内价格拟合而忽略现金流、行权、到期、展期或信用事件，即使误差较小，也可能无法生成与实际损益一致的持有期风险；反之，如果只采用刚性的理论模型而不处理波动聚集、状态变化和横截面异质性，也难以适应真实市场数据。

\textbf{（1）股票价值路径与波动率建模。}
股票估值与风险预测的研究重点，已经由单一收益率均值逐步转向条件波动率、偏度、厚尾和联合分布。Engle 提出的 ARCH 模型以条件异方差描述波动聚集，为 GARCH、随机波动率和多变量波动模型奠定了基础\cite{engle_1982}；Heston 模型进一步允许波动率随机演化并与标的收益相关，为偏斜分布和衍生品一致估值提供了可计算框架\cite{heston_1993}。此类模型能够从同一随机路径导出未来价格、波动率和尾部损失，但对参数稳定性、极端跳跃、停牌及公司行为的处理仍需结合具体市场规则。对于全量股票担保品，单资产估计还会面临短样本和参数噪声问题，因此需要在保留可识别随机机制的同时，引入分层估计、横截面共享和时间外校准。

\textbf{（2）固定收益、基金和信用风险估值。}
债券价值首先取决于票息、本金和应计利息等现金流，再取决于无风险曲线、信用利差、流动性和违约回收。Nelson--Siegel 模型以少量参数描述收益率曲线的水平、斜率和曲率，为债券现金流折现和曲线情景生成提供了经典基准\cite{nelson_siegel_1987}；Duffie 和 Singleton 的约化型框架则将违约强度与回收处理纳入含违约证券的期限结构估值，可延伸至信用利差期权和信用衍生品\cite{duffie_singleton_1999}。基金并不存在脱离产品类型的统一定价公式：开放式基金通常以净值和持仓穿透为基础，交易型产品还需处理市场价格、折溢价、申赎与跟踪误差。因此，债券、基金和信用衍生品可以共享利率、信用和市场因子，但其统一价值仍必须分别绑定现金流、净值口径和合约条款，不能以收益率、折溢价或信用利差本身替代持有期价值。

\textbf{（3）期货、期权及无套利约束。}
Black 对远期、期货和商品期权的研究明确区分了远期合约固定交割价格与期货逐日结算机制，并给出了基于期货价格的期权估值方法\cite{black_1976}；Black--Scholes 模型及其后续随机波动率扩展，则构成场内期权理论价值、隐含波动率和 Greeks 计算的基础\cite{black_scholes_1973,heston_1993}。近年来，神经网络和序列模型被用于拟合期限结构与隐含波动率曲面，但其有效性仍受无套利条件、稀疏报价和到期边界约束。对担保品风险而言，期货的合约乘数、逐日盯市、保证金和展期规则，期权的权利金、行权、到期和波动率曲面演化，都必须进入未来价值和实际 PnL；只预测结算价或隐含波动率的一部分，不能自动构成完整风险计量器。

\textbf{（4）从点估值到概率分布与流动性风险。}
传统估值通常给出一个价格或收益率，风险管理则需要回答该估计的不确定性以及尾部损失。Artzner 等提出一致性风险度量的基本性质，指出仅依赖分位数的 VaR 在一般情形下可能不满足次可加性，并推动了 ES 等尾部均值度量的发展\cite{artzner_1999}。巴塞尔市场风险框架也将压力期 ES、不可建模风险因子和回测要求纳入内部模型方法\cite{bcbs_2019_market_risk}。与此同时，Amihud 的非流动性指标以绝对收益与成交额之比近似价格冲击，为缺少高频报价时的长期流动性研究提供了可操作变量\cite{amihud_2002}。由此可见，波动率、VaR、ES、成交活跃度和产品敏感度是对持有期损失分布的不同刻画，而不是可以相互替代的“模型数量”。

\textbf{（5）跨序列概率预测与现代学习方法。}
随着大规模面板数据和深度学习的发展，研究开始由逐资产拟合转向跨序列共享。DeepAR 通过在大量相关序列上训练自回归循环网络输出未来概率分布\cite{salinas_2020_deepar}；Temporal Fusion Transformer 将静态变量、历史变量和已知未来变量分开编码，并面向多期限预测提供变量选择和注意力解释\cite{lim_2021_tft}。这些方法为六类产品的分位数、区间和多期限预测提供了可借鉴的计算结构，但其原始研究目标并不等于证券担保品估值：它们通常不自动满足债券现金流恒等式、期权无套利、期货结算或信用事件约束，也不能仅凭跨数据集结果证明在当前样本上优于 last-value、理论模型或强统计基准。

现有主要技术路线与本课题需求的关系如表~\ref{tab:review-product-methods} 所示。

\begin{table}[htbp]
  \centering
  \caption{六类证券担保品估值与风险计量技术比较}
  \label{tab:review-product-methods}
  \ReportTableText
  \begin{tabularx}{\textwidth}{@{}p{4.8em}p{9.5em}p{12.5em}X@{}}
    \toprule
    产品 & 经典技术主线 & 现代扩展方向 & 对统一持有期风险的主要限制 \\
    \midrule
    股票 & 条件波动、随机波动率、历史模拟 & 分层/面板学习、概率序列模型、状态条件化 & 公司行为、停牌涨跌停、厚尾和全量横截面稳定性需显式处理 \\
    债券 & 现金流折现、收益率曲线、久期与信用强度 & 曲线/利差残差学习、跨发行人面板模型 & 全价、应计、票息、本金、违约和流动性不能由单一收益率代理 \\
    基金 & 净值、持仓复制、市场因子与跟踪误差 & 持仓穿透、序列/面板残差模型 & 基金子类、净值与市场价、申赎分红和折溢价口径不同 \\
    期货 & 持有成本、基差和期限结构 & 期限结构序列模型、状态与滚动合约学习 & 逐日盯市、乘数、保证金、到期和展期规则决定实际 PnL \\
    期权 & Black--Scholes/Black--76、波动率曲面与 Greeks & 无套利约束曲面、标的与波动率联合预测 & 稀疏报价、静态套利、行权到期和长持有期可评价性 \\
    信用衍生品 & 强度/生存/回收模型与双腿现金流 & 信用面板、语义事件和残差模型 & 合约、报价、信用事件和回收数据缺一不可，代理或合成数据不能替代真实验证 \\
    \bottomrule
  \end{tabularx}
\end{table}

\textbf{（6）中外金融行业实践与制度体系。}
从实际应用看，国内金融机构通常把估值、折算率、保证金和返回检验嵌入分业监管规则、交易场所参数和机构内部限额；国际主要市场则更多同时受巴塞尔框架、PFMI、FSB 证券融资折扣框架以及 ISDA 标准化保证金方法影响。两类体系的共同目标都是把市场价格变化、流动性、信用质量、集中度和处置期损失转化为可执行的担保品价值或保证金要求，但适用对象、法律结构和参数粒度并不相同。表~\ref{tab:review-industry-practice-cn-intl} 从业务使用角度归纳其主要差异。

\begin{table}[htbp]
  \centering
  \caption{中外金融行业证券担保品估值与风险管理实践比较}
  \label{tab:review-industry-practice-cn-intl}
  \ReportTableText
  \begin{tabularx}{\textwidth}{@{}p{6.5em}XXX@{}}
    \toprule
    业务领域 & 国内行业常用做法 & 国际主要市场常用做法 & 比较及本课题借鉴 \\
    \midrule
    商业银行交易账簿 & 按监管标准法或经批准的内部模型法计量市场风险；内部模型与预期尾部损失、损益归因、返回检验、压力测试和独立验证联动\cite{nfra_capital_2023,nfra_advanced_methods_2024} & FRTB 以标准法和内部模型法并行，内部模型强调 ES、不可建模风险因子、损益归因和交易台返回检验\cite{bcbs_2019_market_risk} & 监管框架已较为接近，但监管资本不等于担保品估值；本课题分别保存理论价值、实际 PnL、风险量和规则参数 \\
    融资融券与证券融资 & 采用可充抵证券名单、折算率上限、逐日维持担保比例、集中度约束和差异化动态调整，并以压力测试支持参数调整\cite{csrc_margin_financing_2015,sse_margin_rules_2023} & 以市值折扣、逐日盯市、变动保证金、集中度及错向风险控制为主；FSB/Basel 对部分非集中清算证券融资规定方法标准和最低折扣率\cite{fsb_sft_haircuts_2020,basel_sft_haircut_floors_2021,ecb_collateral_risk_control} & 国内更突出交易场所名单与上限管理，国际框架更突出风险敏感折扣和跨机构可比；模型建议不得突破现行规则边界 \\
    交易所与中央对手清算 & 中央对手对担保品每日估值并设置折扣率，执行日终或日间盯市、保证金、清算基金、回溯测试和压力测试\cite{shclearing_ccp_rules_2026,shclearing_pfmi_2024} & PFMI 要求风险敏感且定期复核的保证金体系；CME SPAN/SPAN 2 以情景风险阵列、组合抵扣和附加项计量期货期权保证金\cite{cpmi_iosco_pfmi_2012,cme_span_methodology} & 两者均形成“估值--保证金--回测--压力”闭环；本课题采用该闭环，但只在有证据时确认跨品种抵扣 \\
    场外金融衍生品 & 人民币利率互换等业务实行集中清算，双边业务以 NAFMII 主协议及履约保障文件约定估值、净额结算和变动保证金\cite{pbc_otc_clearing_2014,nafmii_vm_2025} & 非集中清算衍生品普遍按 BCBS--IOSCO 框架交换初始/变动保证金，初始保证金常以 ISDA SIMM 统一敏感度口径\cite{bcbs_iosco_margin_2015,isda_simm_2025} & 必须区分集中与双边清算、初始与变动保证金，并显式记录阈值、净额组、担保品法律形式和处置期 \\
    模型验证与运行 & 监管和清算机构强调独立验证、返回检验、损益归因、压力测试、数据质量以及模型持续应用\cite{nfra_advanced_methods_2024,shclearing_pfmi_2024} & PFMI 与行业模型治理强调日常回测、敏感性分析、定期校准、版本治理和定量披露\cite{cpmi_iosco_pfmi_2012,isda_simm_2025} & 国内外均不接受“一次建模、长期不变”；本课题需保留版本、数据快照、参数、告警和回退证据 \\
    \bottomrule
  \end{tabularx}
\end{table}

表~\ref{tab:review-standards-cn-intl} 进一步按制度功能对照国内规则与国际标准。这里的“对照”用于识别共同控制目标和技术接口，不表示法律等同、互认或逐条一致。

\begin{table}[htbp]
  \centering
  \caption{中外证券担保品估值与风险管理相关标准、监管规则及行业协议对照}
  \label{tab:review-standards-cn-intl}
  \ReportTableText
  \begin{tabularx}{\textwidth}{@{}p{6.5em}XXX@{}}
    \toprule
    制度功能 & 中国标准或规则 & 国际标准或主要参照 & 对照关系与适用边界 \\
    \midrule
    公允价值计量 & 《企业会计准则第39号——公允价值计量》规定有序交易中的退出价格，并规范估值技术与公允价值层次\cite{cas39_fair_value_2014} & IFRS 13 统一公允价值定义、计量框架和披露要求\cite{ifrs13_fair_value_2011} & 核心定义和层次具有较强可比性；具体适用主体、披露和其他准则衔接仍按各自会计体系执行 \\
    银行市场风险资本 & 《商业银行资本管理办法》及高级方法验收规定覆盖标准法、内部模型法、ES、损益归因、返回检验和持续监管\cite{nfra_capital_2023,nfra_advanced_methods_2024} & 巴塞尔市场风险最低资本要求（FRTB）规定标准法、内部模型法及交易台级验证\cite{bcbs_2019_market_risk} & 国内规则吸收了 FRTB 的主要技术结构，但银行分档、监管审批和实施范围以中国规则为准 \\
    证券融资折扣 & 证监会和交易所规则要求对可充抵证券、折算率、维持担保比例和集中度进行动态、差异化管理\cite{csrc_margin_financing_2015,sse_margin_rules_2023} & FSB 证券融资折扣框架与 Basel CRE56 对特定非集中清算交易提出定性标准和最低折扣率\cite{fsb_sft_haircuts_2020,basel_sft_haircut_floors_2021} & 均用于抑制杠杆与顺周期性，但交易范围、豁免对象、资产分类和数值表并非一一对应 \\
    金融基础设施与中央清算 & 2013 年起以 PFMI 评估境内基础设施，2025 年《金融基础设施监督管理办法》进一步要求与 PFMI 等国际准则接轨\cite{pbc_csrc_pfmi_assessment_2014,pbc_csrc_fmi_rules_2025} & CPMI--IOSCO《金融市场基础设施原则》规定担保品、保证金、信用风险、流动性和违约管理原则\cite{cpmi_iosco_pfmi_2012} & 属于明确的原则衔接；准入、监管权限、业务规则和具体参数仍由境内制度及基础设施确定 \\
    非集中清算衍生品保证金 & NAFMII 主协议和 2025 年变动保证金履约保障文件提供境内银行间市场的标准合同安排\cite{nafmii_vm_2025} & BCBS--IOSCO 规定初始/变动保证金最低框架，ISDA SIMM 提供行业统一初始保证金模型\cite{bcbs_iosco_margin_2015,isda_simm_2025} & 控制目标相近，但适用门槛、合格担保品、净额结算、隔离和法律实现方式需逐交易核验 \\
    \bottomrule
  \end{tabularx}
\end{table}

上述比较表明，国内与国际实践在公允价值、逐日盯市、风险敏感折扣、压力测试和模型验证等原则上具有较高共性，差异主要集中在制度适用范围、折扣参数的制定主体、组合抵扣、法律可执行性和数据披露粒度。因此，本课题应在系统中同时保留两类字段：一类是来自监管、交易所、清算机构或合同的刚性规则参数，另一类是模型估计的未来价值、流动性冲击和尾部损失。业务担保品价值和风险预警应在适用规则下调用模型结果，而不能用模型输出覆盖法定折算率、保证金下限或合同约定。

总体而言，现有研究已经为六类产品分别提供了较成熟的定价与风险组件，但仍存在三方面断点：一是研究对象经常在价格、收益率、利差、波动率和流动性指标之间切换，缺少统一的价值与损失语义；二是理论模型与机器学习模型常使用不同样本、不同信息集和不同评价指标，横向结论缺乏可比性；三是估值、持有期 PnL、VaR/ES 和系统监控常由不同流程生成，难以形成可追溯证据。本课题据此不采用“一套算法覆盖所有产品”的方式，而是在共同输出契约下建立 M1--M6 六个产品主模型，使每类模型都从同一条件价值路径产生 \(V_t\)、\(V_{t+h}\)、期内现金流、PnL、正损失分布、VaR 和 ES，再以产品特有指标解释差异。

\section{大模型金融时序增强技术的研究进展}

文本与大模型在金融预测中的应用，可以概括为“词典与情绪指标—领域预训练编码器—金融大语言模型—语言和时间序列联合建模”四个阶段。该领域的共同出发点是：公告、政策、新闻、产品说明和主体披露包含数值序列未完全反映的状态信息；其共同难点则是发布时间、语义噪声、样本选择和模型容量可能制造虚假增益。对于证券担保品估值，关键问题并不是大模型能否生成一段合理的金融解释，而是当预测时点只能使用当时已发布的信息时，语义变量能否在同一底座、同一样本和同一评价口径下稳定改善未来价值或损失分布。

\textbf{（1）金融文本指标与领域语义编码。}
早期研究主要将新闻或披露文本转化为可进入计量模型的情绪和风险变量。Tetlock 以媒体内容量化投资者情绪并检验其与市场价格和成交量的关系\cite{tetlock_2007}；Loughran 和 McDonald 指出通用负面词典会误判金融语境中的大量词语，说明金融文本必须进行领域化处理\cite{loughran_mcdonald_2011}。在预训练语言模型阶段，FinBERT 通过金融语料适配 BERT，用较少标注数据改善金融情绪分类\cite{araci_2019_finbert}。这些工作证明了文本可以形成结构化条件变量，但“文本分类准确”与“估值或尾部风险预测增益”仍是两个不同问题，后者必须用未来市场实现值另行检验。

政策语义是金融文本状态中的重要组成。Baker、Bloom 和 Davis 通过新闻报道频率构建经济政策不确定性指数，并检验其与波动、投资及宏观活动的关系\cite{baker_2016_epu}。该类指数为研究政策不确定性与资产风险的关联提供了可复现入口，但它仍属于特征而不是预测模型；若指数发布、修订或匹配时间晚于 \texttt{as\_of} 时点，或者只在事后选择有利滞后期，其表面增益可能来自时间错配而非真实可用信息。

\textbf{（2）金融大语言模型与检索增强。}
BloombergGPT 采用金融与通用语料混合训练，显示领域数据能够提升多类金融自然语言任务，同时保留一定通用能力\cite{wu_2023_bloomberggpt}。检索增强生成（RAG）则把参数化语言模型与外部非参数知识库结合，用检索结果为生成内容提供更新入口和来源线索\cite{lewis_2020_rag}。对担保品场景而言，领域大模型和 RAG 更适合承担文本筛选、事件抽取、条款解析、状态编码和证据定位，而不宜直接覆盖数值定价模型。原因在于生成式答案即使语言流畅，也不天然满足现金流恒等式、无套利边界和损失正负号约定；只有把检索文档、发布时间、模型版本和数值残差同时保存，才能审计语义修正的来源。

\textbf{（3）大模型与时间序列基础模型。}
Time-LLM 通过时间序列重编程和提示前缀，将数值片段映射到冻结语言模型的表示空间，探索了少样本和零样本预测\cite{jin_2024_timellm}；Chronos 则对数值进行缩放、量化和离散化，使用语言模型架构训练概率时间序列模型\cite{ansari_2024_chronos}。这些研究推动了跨数据集预训练和统一预测接口，但也引出了“大语言模型中的语言预训练究竟贡献了什么”的识别问题。Tan 等对多种 LLM 时间序列方法进行消融后发现，在其评测范围内，移除预训练语言模型或替换为基础注意力层并未导致性能下降\cite{tan_2024_llm_ts}。这一负向证据说明，模型规模、参数冻结和语言模型名称不能代替同样本增量检验；数值 patch、注意力结构和训练数据本身可能已经解释了大部分性能。

\textbf{（4）受约束残差增强与可识别增量。}
综合上述研究，金融大模型的可行定位应从“直接生成价格”转向“在产品底座上提供受约束的语义残差或状态门控”。一方面，底座模型负责现金流、产品条款、无套利边界和基础损失分布；另一方面，语义模块只使用预测时点已发布的政策、公告、新闻和产品文本，输出幅度受限的残差、分位数修正或状态权重。当文本缺失、冲突、延迟、提示异常或输出越界时，系统确定性回退到底座。更重要的是，真实文本版本必须同时优于同架构 \texttt{no\_text}、打乱文本和时间错位文本，才能把差异解释为语义增量。

表~\ref{tab:review-llm-methods} 对相关技术路线及其在本课题中的定位进行了归纳。

\begin{table}[htbp]
  \centering
  \caption{金融文本与大模型时序增强技术比较}
  \label{tab:review-llm-methods}
  \ReportTableText
  \begin{tabularx}{\textwidth}{@{}p{7.5em}p{10.5em}p{10.5em}X@{}}
    \toprule
    技术路线 & 主要输入与输出 & 优势 & 主要风险及本课题定位 \\
    \midrule
    金融词典/文本指标 & 新闻、公告或披露文本；输出情绪、语调和主题指标 & 计算透明，便于时间序列对齐和经济解释 & 词义依赖领域且覆盖有限；作为可审计特征和简单基准 \\
    领域预训练编码器 & 金融语料与任务标注；输出文本向量或分类概率 & 少量标注下可迁移，语义表达强于固定词典 & 分类性能不等于风险增益；作为 L1/L2 的编码器候选 \\
    金融大语言模型/RAG & 领域语料、检索文档和提示；输出事件、条款、摘要或证据 & 能处理长文本与多任务，并可绑定外部知识 & 幻觉、提示注入和来源漂移；限于证据抽取、状态编码与解释 \\
    LLM 时间序列重编程 & 数值 patch、提示和冻结语言模型；输出点或分布预测 & 统一接口，具有少样本迁移潜力 & 语言预训练贡献存在争议；必须与去除 LLM 的消融同样本比较 \\
    时间序列基础模型 & 大规模跨域数值序列；输出零样本或微调预测分布 & 跨序列共享强，可提供现代挑战基准 & 不自动满足金融产品约束；只作强基准或残差组件 \\
    受约束语义残差 & 产品底座、真实时点文本与质量门禁；输出残差或状态权重 & 保留数值底座，可消融、可回退、可追溯 & 需要严格发布时间、幅度/无套利约束和负向结果披露；为本课题主线 \\
    \bottomrule
  \end{tabularx}
\end{table}

总体而言，金融文本研究已经证明领域语义具有潜在信息价值，基础模型研究也展示了跨任务迁移能力，但从“语言任务有效”到“证券担保品估值及价差风险有效”之间仍缺少三个环节：一是真实发布时间和数据 vintage 约束，二是与产品数值底座的经济边界，三是能够排除模型容量、时间标签和样本差异的对照试验。本课题据此设置 L1 和 L2 两个增强器：L1 检验真实时点文本能否改善估值点值与分位数，L2 检验政策不确定性及其他语义状态能否改善未来价值、损失分位数、波动率和 VaR。两个增强器均保留底座、检索证据、消融和回退，不以生成式表述代替数值验证。

\section{返回检验与动态监控技术的研究进展}

估值及风险模型只有在预测完成后与真实实现值比较，才能形成可证伪的科学结论。返回检验技术早期集中于 VaR 例外次数，随后扩展到例外独立性、完整预测分布、VaR/ES 联合评价和模型间预测比较；动态监控则进一步关注输入分布、预测误差、覆盖率和系统状态随时间的变化。现有研究的基本共识是，单一准确率无法同时回答“点预测是否精确、分布是否校准、尾部风险是否可靠、异常是否聚集以及系统是否仍在有效运行”等问题，必须建立分层指标和连续证据链。

\textbf{（1）VaR 例外覆盖与独立性检验。}
Kupiec 提出的比例失败检验以实际损失超过 VaR 的次数检验无条件覆盖率，并指出低概率尾部事件在小样本下检验能力有限\cite{kupiec_1995}。Christoffersen 进一步把正确覆盖与例外独立性结合，形成条件覆盖评价框架：合格的风险区间不仅应具有正确的总体命中比例，例外也不应在时间上异常聚集\cite{christoffersen_1998}。巴塞尔市场风险规则以实际或假设 PnL 与风险度量比较，并按例外情况采取分层监管响应\cite{bcbs_2019_market_risk}。因此，经验覆盖率达到某一阈值只是必要信息，不能单独证明模型在压力期、不同产品或不同持有期上均有效。

\textbf{（2）概率预测与 VaR/ES 联合评价。}
Gneiting 和 Raftery 系统论述了严格 proper scoring rule，说明概率预测应使用能够鼓励真实分布报告的评分规则；其中 CRPS、分位数损失和区间评分分别适用于完整分布、分位数和预测区间\cite{gneiting_raftery_2007}。Fissler 和 Ziegel 证明在适当条件下 VaR 与 ES 可以联合诱导，从而为二者的联合预测比较提供一致评分基础\cite{fissler_ziegel_2016}。这意味着模型评价不能只检查 VaR 是否“少超限”，还应比较预测分布的尖锐性与校准性、ES 对尾部严重程度的刻画以及点估值误差。过宽的风险区间可能获得较高覆盖率，却缺少决策价值；只优化均方误差的模型也可能严重低估尾部损失。

\textbf{（3）同样本预测比较与时间依赖。}
Diebold--Mariano 检验以损失差序列比较两个预测的准确性，允许使用非对称、非二次损失并处理一定的序列相关\cite{diebold_mariano_1995}。对于多期限金融预测，相邻起点的 \(h\) 日损益窗口会机械重叠，若把这些观测视为独立样本，置信区间和例外聚集判断可能失真。Newey--West 的异方差自相关一致协方差估计\cite{newey_west_1987}和 K{\"u}nsch 的块自助法\cite{kunsch_1989}为相关序列的不确定性估计提供了基础，但 VaR 独立性检验仍需在非重叠序列上实施。因而，多期限评价应预先区分用于误差估计的全量重叠样本和用于独立性判断的非重叠 cohort，不能在看到结果后选择最有利的期限或抽样方式。

\textbf{（4）数据漂移、性能漂移与在线监控。}
模型上线后的风险不仅来自参数误差，也来自输入分布、产品构成、数据质量和市场机制变化。Bifet 和 Gavald{\`a} 提出的 ADWIN 根据数据流变化自适应调整窗口，可用于监测误差率或条件概率漂移\cite{bifet_gavalda_2007}；Gretton 等提出的最大均值差异（MMD）以核均值嵌入比较两个样本分布，为多变量输入或残差分布变化提供了非参数检验工具\cite{gretton_2012}。然而，分布变化并不必然意味着模型失效，模型误差上升也可能来自数据缺失、条款变更或覆盖分母变化。因此，动态监控需要同时观察输入、输出、实现 PnL、例外、缺失、回退和 catalogue 构成，并通过异常诊断区分数据问题、模型问题与真实市场状态变化。

\textbf{（5）从统计检验到可审计证据链。}
现有返回检验研究主要回答统计意义上的校准问题，工程系统还需保证每个结果可被重建。预测必须绑定当时可用的信息截止时点、数据快照、模型版本、配置和代码哈希；历史回放必须标明它是事后重建还是当日真实留痕；停牌、到期、违约、无成交和缺失数据必须按照事前规则进入排除台账。只有将预测明细、实现值、例外标志、指标、日志和版本 manifest 联结起来，才能避免手工复制表格或事后修改分母造成报告、平台和测试材料之间的分叉。

表~\ref{tab:review-validation-methods} 汇总了返回检验和监控方法的适用问题及边界。

\begin{table}[htbp]
  \centering
  \caption{返回检验与动态监控方法比较}
  \label{tab:review-validation-methods}
  \ReportTableText
  \begin{tabularx}{\textwidth}{@{}p{8em}p{10em}p{10em}X@{}}
    \toprule
    验证问题 & 代表方法 & 能够回答的问题 & 主要边界及本课题处理 \\
    \midrule
    点值与分位数精度 & MAE/RMSE、WAPE/sMAPE、pinball loss & 当前/未来价值及损失分位数误差大小 & 不代表完整分布或合同通过；按产品、期限、市况和流动性分层 \\
    完整概率分布 & CRPS、区间覆盖与宽度、严格 proper score & 分布是否兼顾校准与尖锐性 & 需相同样本和同一损失基础；与点值、波动率指标联合报告 \\
    VaR 覆盖与聚集 & Kupiec、Christoffersen、例外时间图 & 例外比例是否正确、是否存在异常聚集 & 小样本检验能力有限；多日持有期在非重叠 cohort 内检验 \\
    VaR/ES 尾部评价 & 联合一致评分、压力期分层 & 尾部分位数与尾部均值是否协调 & 不能只以高覆盖掩盖过宽风险值；同时报告严重度和压力期 \\
    模型间预测比较 & 配对损失差、Diebold--Mariano、块 bootstrap/HAC & 同一样本上增益是否具有统计和经济意义 & 嵌套模型、多期限和多重比较需预注册并修正 \\
    动态漂移监控 & ADWIN、MMD、滚动误差与覆盖率 & 输入、残差或性能是否随时间变化 & 告警不等于模型失效；结合数据质量、回退和市场状态诊断 \\
    工程与覆盖审计 & catalogue、预测/排除台账、版本与哈希 manifest & 全品种分母、失败原因和结果是否可追溯 & 属于工程与覆盖证据，不能由统计精度自动替代 \\
    \bottomrule
  \end{tabularx}
\end{table}

总体而言，现有方法已经建立了从点预测、概率预测到 VaR/ES 的较完整评价体系，也提供了漂移检测和在线更新工具，但把这些方法用于六类证券担保品仍需解决三项衔接问题：第一，风险预测与实现 PnL 必须采用一致的价值、现金流和损失正负号；第二，五年历史回放、多持有期窗口和全品种覆盖会引入重叠相关、样本资格变化和缺失机制；第三，统计通过、科学增益、工程稳定和产品覆盖回答的是不同问题，不能合并成单一“准确率”。本课题因此把返回检验与动态监控设计为独立研究点，在统一预测台账上同时执行估值误差、概率评分、波动率评价、VaR/ES 检验、漂移监测和证据归档。

\section{本章小结}

通过对多品种估值、大模型时序增强以及返回检验与动态监控技术的综述，可以形成以下认识。

第一，六类证券担保品具有不同的定价机制和现金流事件，股票波动率模型、债券曲线与信用模型、基金净值与持仓方法、期货期限结构、期权无套利模型以及信用衍生品强度模型均有成熟研究基础，但这些基础尚不能自然汇合为统一的跨产品风险结果。可行路径是在产品内部保留经济约束，在产品之间统一 \(V_t\)、\(V_{t+h}\)、现金流、PnL、正损失、VaR 和 ES 的输出语义。

第二，深度概率预测和时间序列基础模型能够共享横截面信息并输出多期限分布，但不能脱离产品机制直接替代经典底座。现代模型是否有效，必须在当前数据、相同信息集和相同 eligible sample 上与 last-value、产品经典模型及强现代基准比较，不能引用不可比论文数据宣布优越性。

第三，金融文本、领域语言模型和检索增强为政策、公告、新闻与条款信息的利用提供了工具，但语义分类能力并不等于估值与风险增量。大模型在本课题中应作为受约束、可消融、可回退的残差或状态增强器；真实文本版本只有在优于无文本、打乱文本和时间错位文本，并且不恶化合同风险指标时，才能被认定为产生增量技能。

第四，返回检验已经从例外次数扩展到覆盖独立性、概率评分、VaR/ES 联合评价和模型间配对比较，动态监控也能够检测输入与性能漂移。但针对多产品、多期限和五年历史回放，仍需统一现金流与损失基础、处理重叠窗口、冻结 catalogue 分母，并把统计结果与数据、代码、配置和日志建立可审计关联。

基于上述研究进展和不足，第 3 章将研究 M1--M6 六类产品主模型及其统一持有期损失输出；第 4 章将在产品底座之上研究 L1/L2 语义增强、消融、约束和回退机制；第 5 章将研究五年实际市场返回检验、动态监控、异常诊断和证据归档。三章分别回答“如何计量”“语义是否真正增益”和“结果能否由真实市场及持续运行验证”，共同构成从估值到风险、从方法到证据的闭环。
